\begin{textsample}{NEG Dim 4 – pi – Score -1.00 – pi/pi\_em\_33.txt}  \label{ex:f4_neg_244}
Educação Física Revisitando a trajetória da Educação Física nas décadas anteriores aos anos 80, percebe -se que além das próprias limitações, condições de trabalho, valorização do profissional, na prática este componente curricular ainda apresenta marcas de um modelo biológico de aptidão, utilitário.
Diferentes pressupostos, muitas vezes contraditórios e antagônicos, têm sido defendidos e difundidos através do discurso e da prática pedagógica dos profissionais e estudiosos dessa área.
Na década de setenta, a luta dos profissionais e estudiosos da Educação Física tinha a finalidade de inserir e fazer cumprir de fato a obrigatoriedade desse componente no âmbito escolar resultando, nas reformas de 1971, em que a educação Física tornou -se obrigatória em todos os níveis de escolarização ; a Lei de Diretrizes e Base – LDB ( 1996 ), em que passou a componente curricular com a mesma importância que todos os outros, com novos arranjos relacionadas a estrutura didática e autonomia das escolas e sistemas de ensino.
Os PCNs ( 1999 ) apontam propostas para desenvolvimento do componente, orientando professores a trabalhar de forma lúdica sem perder o caráter educativo.
E ainda, coloca a Educação Física como área de estudo fundamental para compreensão do ser humano, enquanto produtor de cultura.
Atualmente, o desafio desses profissionais encontra -se na redefinição da identidade da educação física escolar, considerando que a relação objetivo-conteúdo – forma suporte a re-significação do conhecimento, como prática social e cultural integrada com a \textbf{consecução} das finalidades da instituição escolar.
Esse desafio implica em superar a histórica da subordinação da identidade da educação física escolar às tendências, aos valores, aos códigos e às orientações conceituais e metodológicas que têm sido forjadas a partir de qualquer instituição externa à escola como assinala Bracht, ( 1991 ).
As produções intelectuais que no início tinham preocupação com a dimensão da legalidade, agora discorrem sobre a questão da legitimidade da Educação Física.
Esse é um dos desafios premente, na medida em que as conjecturas pedagógicas projetadas para o ensino da Educação Física, não são suficientemente capazes de dar conta da legitimação desse componente curricular no âmbito escolar na atualidade.
Estas questões enseja algumas considerações cabendo reconhecer, que a realidade é bem mais complexa do que os esquemas mentais para explicá -la, portanto neste momento de atualização dos documentos curriculares, não podemos deixar passar despercebido as questões culturais, os universos simbólicos e as inquietações das crianças, jovens no sentido compreender e atender os sujeitos dessa etapa da Educação Básica, uma vez que temos um quadro forte de evasão e abandono desses jovens.
Neste ponto Carrano ( 2013, p. 15 ) destaca que a juventude vem ganhando : > [ … ] contornos próprios em contextos históricos, sociais e culturais distintos.
As distintas condições sociais ( origem de classe e cor da pele, por exemplo ), a diversidade cultural ( as identidades culturais e religiosas, os diferentes valores familiares, etc. ), a diversidade de gênero ( a heterossexualidade, a homossexualidade, a transexualidade ) e até mesmo as diferenças territoriais se articulam para a constituição das diferentes modalidades de se vivenciar a juventude.
Além das marcas da diversidade cultural e das desiguais condições de acesso aos bens econômicos, educacionais e culturais, a juventude é uma categoria dinâmica.
Ela é transformada no contexto das mutações sociais que vêm ocorrendo ao longo da história.
Na realidade, não há tanto uma juventude e sim jovens, enquanto sujeitos que a experimentam e a sentem segundo determinado contexto sociocultural em que se inserem e, assim, elaboram determinados modos de ser jovem.
É nesse sentido que enfatizamos a noção de juventudes, no plural, para enfatizar a diversidade de modos de ser jovem existente.
Por conseguinte, quando se fala em estudantes do Ensino Médio deve -se fazer referência a sujeitos das diferentes juventudes que estão em processo de construção identitária.
A solução dessa tensão pode -se dar considerando o critério último que justifica essas posições : o ser humano.
Surge a pergunta, a Educação Física, hoje como está sendo desenvolvido nas escolas promove a dignidade humana, possibilitando -lhes expressar e desenvolver todos os seus talentos?
Conhecimento, experiências?
Promove relações harmônicas entre pessoas, entre grupos sociais?
Promove uma relação com a natureza que a reconheço como dádiva que devemos às futuras gerações?
Promove a produção, a circulação e o consumo de bens que assegurem a todos uma vida digna e respeitem o ambiente?
Enfim, promove e inspira -se nela a relação com o transcendente?
O ponto a resolver esta na tensão, entre a memória e o sonho, é a busca pelo equilíbrio que não é meio termo, mas contemplação dessas dimensões numa relação dialógica e dinâmica a partir de discussões como, por exemplo, qual tipo de conhecimento é valorizado e por quem, tendo em vista, que uma posição pode mudar de acordo com o contexto social e histórico.
Este cenário, provoca em toda comunidade escolar uma pausa, para a identificação e análise crítica das concepções auto definidas de abordagens de ensino, examinando -as partir de que contextos estão sendo elaboradas e se justificam, e que referencial didático-pedagógico está viabilizando a síntese teoria-prática.
Nota -se que as tendências ( higienista, militar, pedagogicista, competitivista e popular ) são classificações que caracterizam períodos históricos da Educação Física, esquemas mentais que ajudam na compreensão da realidade.
Porém a realidade é bem mais complexa e estas valorizações vão sendo corrompidas pelas pressões provocadas pelos diferentes estados sucessivos da sociedade.
E mas, estas mudanças não vão acontecendo isoladamente, estão ligadas em relações complexas, requerendo conceitos que lidem com esta complexidade.
Portanto, experimentar e analisar as diferentes formas de expressão que não se alicerçam apenas nessa racionalidade é uma das potencialidades desse componente na Educação Básica.
O documento, Base Nacional Curricular Comum – BNCC ( 2018 ) com o compromisso de aproximação dos aspectos concretos da aprendizagem, reconhece que a Educação Física proporciona densa possibilidade de experiências para crianças, jovens e adultos na educação básica, implicando na expansão e acesso a um extenso universo cultural.
Acrescenta, pois, que o referido componente inclui saberes corporais, experiências estéticas, afetivas, lúdicas e agonista, como proposições que guiam a busca pela qualidade do trabalho pedagógico sem reduzi -lo à racionalidade típica dos saberes científico que, freqüentemente, orienta as práticas pedagógicas.
Nessa dinâmica, a prática educativa deve ser compreendida a partir de um conjunto de conhecimentos organizados em torno de intencionalidades, procurando criar mecanismos de ação que implicam valores, orientações, códigos e semânticas no sentido de tornar exeqüíveis as propostas consentidas no sentido da manutenção, modernização ou transformação cultural.
Portanto, uma prática pedagógica que incorpore a reflexão contínua e coletiva, de forma a assegurar que a intencionalidade proposta seja disponibilizada a todos, como argumenta, Franco ( 2016 ).
Portanto a direção da prática pedagógica da educação física, deve estar alinhada ao princípio da mediação crítica entre teoria e prática, assegurando que os alunos ( re ) significar os conhecimentos, como ferramentas mentais que auxiliem na compreensão a respeito dos seus movimentos e dos recursos para o cuidado de si e dos outros.
E ainda, Desenvolver autonomia para apropriação e utilização da cultura corporal de movimento em diversas finalidades humanas, favorecendo sua participação de forma confiante e autoral na sociedade.
A escola, entendida como “ locus ” político-social de promoção e circulação dos conhecimentos culturais apreendidos e elaborados historicamente, com a função básica de desenvolver habilidades, atitudes e pautas de comportamento para a formação do cidadão para sua intervenção na vida pública, de forma crítica, deve também ser um espaço privilegiado para a prática das atividades corporais.
O planejamento das aulas deve assegurar que as práticas corporais sejam apresentadas como fenômeno cultural dinâmico, diversificado, pluridimensional, singular e contraditório, e nunca como possibilidades de aprendizagem “ espontaneísta ” e automotivada.
Diante do exposto, o ponto primordial para as mudanças passa, obrigatoriamente, por uma atualização do currículo no que cerca a Educação Física.
Currículo como defende Gómez ( 1998 ), fundado no pensamento de ( STENHOUSE, 1984 ) como um documento vivo em que ofereça modelos propícios de compartilhamento de princípios e propostas educativas, com características de flexibilidade, aberto à reflexão crítica e passíveis de serem efetivamente colocados em prática.
Por esta razão, as questões culturais não podem passar despercebidas neste processo de reconstrução do currículo pelos educadores sob o risco da escola se distanciar dos universos simbólicos e das inquietudes das crianças e jovens.
Este Reconhecimento do universo multicultural pelo aluno prepara -o para enfrentamento e discernimento na escolha de respostas efetivas às questões postas pela sociedade, em particular sobre as desigualdades impostas aos sujeitos de forma por vezes inevitável.
Incluem -se neste contexto a busca de práticas sociais/culturais que induzem a solidariedade, a colaboração, a crítica e a criação evocado pela educação multicultural têm influenciado os vários componentes curriculares, dentre eles a educação Física.
A experiência efetiva das práticas corporais deve oportunizar aos alunos participar de forma autônoma, contextos de lazer e saúde, compreendendo a complexidade do corpo humano nos seus diversos aspectos de existência.
Cada prática corporal deve propiciar ao sujeito o acesso a uma dimensão do conhecimento e de experiência aos quais ele não teria de outro modo, como conferir o documento da BNCC ( 2018 ).
Conclui -se que a vivência da prática corporal suscita conhecimentos muito específicos, expressivos, porém, necessita como o próprio documento destaca, problematizar, desnaturalizar e vivenciar a multiplicidade de sentidos e significados que os grupos sociais atribuem às diferentes manifestações da cultura corporal de movimento.
A Educação Física escolar, neste documento, como atividade curricular versa pedagogicamente, da reflexão e da prática de conhecimentos e habilidades do que estamos denominando de cultura corporal.
Esta se configura através de temas ou formas de atividades particularmente corporais, tais como : os jogos recreativos e esportivos ; a ginástica ; a dança ; as manifestações folclóricas, dentre outras práticas que constituem seu conteúdo.

% matched lemmas: consecução
\end{textsample}
